% ===================================================================
% Table: Comparison with state-of-the-art methods
% ===================================================================
\begin{table}[htbp]
  \centering
  \caption{与同期骨骼手势识别方法的对比}
  \label{tab:cross_method_comparison}
  \small
  \setlength{\tabcolsep}{8pt}
  \begin{tabular}{lcccccc}
  \hline
  方法 & 年份 & 参数量 & 时延 (ms) & 类别 & 平台 & F1 \\
  \hline
  ST-GCN \cite{yan2018stgcn} & 2018 & 3.10M & ${>}$10 & 14/28 & GPU & 0.927$^*$ \\
  DD-Net \cite{yang2019ddnet} & 2019 & 0.15M & ${<}$1 & 14/28 & GPU & 0.946$^*$ \\
  HAN \cite{han2025hierarchical_attention} & 2025 & 0.05M & ${<}$1 & 14/28 & GPU & 0.951$^*$ \\
  Habib et al.\ \cite{habib2025skeleton_hgr} & 2025 & 0.12M & ${<}$1 & 14/28 & CPU & 0.943 \\
  Zhang et al.\ \cite{openxr2025gesture} & 2025 & 0.42M & 1.3 & 10 & XR2 & 0.950 \\
  \hline
  \textbf{DBEW-NN (本文)} & 2026 & \textbf{0.057M} & \textbf{1.1} & \textbf{11} & \textbf{XR2} & \textbf{0.889} \\
  \hline
  \end{tabular}
  \noindent\begin{minipage}{\linewidth}\footnotesize
  注：$^*$为公开基准（SHREC'17）结果，非端侧实测，仅供量级参考。本文方法参数量最小（57K），所有指标来自骁龙XR2端侧实测，含完整ONNX$\rightarrow$Barracuda部署与鲁棒性评估。
  \end{minipage}
\end{table}
